\documentclass[preprint,12pt]{elsarticle}

%% Essential packages for math, figures, and algorithms
\usepackage{amssymb}
\usepackage{amsmath}
\usepackage{graphicx}
\usepackage{subcaption}
\usepackage{booktabs}    %% For professional tables
\usepackage{algorithm}   %% For algorithm environment
\usepackage{algorithmic} %% For algorithm pseudocode
\usepackage[table]{xcolor}
\usepackage[colorlinks=true, citecolor=blue, linkcolor=black]{hyperref}

%% Journal metadata (placeholder – update as needed)
\journal{Computers \& Industrial Engineering}
\biboptions{sort&compress}

\begin{document}

\raggedbottom

\begin{frontmatter}

\title{A Reinforcement Learning and Machine Learning Framework for Adaptive Optimization of the Vehicle Routing Problem with Synchronized Visits}

%% Author list (placeholder – update with correct information)
\author[1]{Nada Farah Djedjig}
\ead{nadadjedjig@gmail.com}
\author[1]{Supervisor Name}
\ead{supervisor@gmail.com}
\author[1]{Supervisor Name}
\ead{supervisor@gmail.com}

\address[1]{Department of Computer Science, University Name, City, Country}

\begin{abstract}
%% ----------------------------------------------------------------------
%% ABSTRACT (150-250 words)
%% ----------------------------------------------------------------------
% Content to write: A concise summary of the entire paper.
% - State the problem: VRPTW-S for home healthcare.
% - State the gap: Metaheuristics like GVNS rely on static decisions, limiting their performance.
% - State the contribution: A novel multi-pronged ML/RL framework to overcome these limitations.
% - State the key results: The framework scales to 400+ customers, DQN is robust, hybrids with RF and PPO provide additional gains.
% ----------------------------------------------------------------------
The Vehicle Routing Problem with Time Windows and Synchronization constraints (VRPTW‑S) is a core optimization challenge in home healthcare logistics, essential for ensuring timely and coordinated care. While state‑of‑the‑art metaheuristics like Generalized Variable Neighborhood Search (GVNS) achieve good performance, their effectiveness is fundamentally limited by three static decisions: (i) a fixed order of neighborhood operators, (ii) a fixed local search mode (First vs. Best Improvement), and (iii) hand‑tuned hyperparameters. To overcome these limitations, this paper introduces a comprehensive, learning-based framework that adaptively guides the GVNS search process. We first enhance the GVNS with eight neighborhood operators and generate a large set of Solomon-based VRPTW‑S instances. Our core contributions are: (1) an RL-based agent for adaptive operator selection, implemented and compared using both tabular Q‑learning and a Deep Q‑Network; (2) a two‑head Random Forest model that dynamically selects the optimal local search mode based on the current search context; and (3) a Proximal Policy Optimization (PPO) controller that manages key hyperparameters online. Extensive experiments demonstrate that DQN scales effectively to 400+ customers, achieving consistent improvements, while a hybrid DQN+RF model yields further gains of up to 2.3\% on the largest instances. The results confirm that integrating machine learning into the optimization loop leads to a more robust, scalable, and high-performing system, providing a promising direction for data-driven logistics.
\end{abstract}

\begin{keyword}
Vehicle Routing Problem \sep Time Windows \sep Synchronization \sep Home Healthcare \sep Reinforcement Learning \sep Deep Q-Network \sep Proximal Policy Optimization \sep Random Forest
\end{keyword}

\end{frontmatter}

% ----------------------------------------------------------------------
%% SECTION 1: INTRODUCTION
% ----------------------------------------------------------------------
\section{Introduction}
\label{sec:intro}

%% ----------------------------------------------------------------------
% Content to include (from report Chapters 1 & 4):
% 1. The motivation: Growing home healthcare demand → need for efficient logistics.
%    - Source: Report, Chap. 1 §1.1 (Home health care routing and Optimization problems).
%    - Key points: Aging population, chronic diseases, cost containment → HHC expansion.
% 
% 2. The problem: VRPTW-S as a hard combinatorial optimization problem.
%    - Source: Report, Chap. 1 §1.2 & §1.3 (VRP & VRPTW & VRPTW-S).
%    - Key points: NP-hard, synchronized visits create tight interdependencies.
%    - Complexity: Exact methods fail for realistic instances (n > 20 customers).
%    - Real-world relevance: Home healthcare routing and scheduling.
%
% 3. The gap: The state of the art (metaheuristics like GVNS) has three static 
%      decisions that limit performance.
%    - Source: Report, Chap. 4 §4.5 (Improvement opportunities for ML/RL).
%    - Key points: (i) Fixed operator order, (ii) Fixed choice between First and Best 
%      Improvement, (iii) Hyperparameters tuned once offline (DoE).
%
% 4. The solution: A multi-part learning-based framework to adaptively guide the search.
%    - Briefly introduce the three main contributions: RL for operator selection, ML for 
%      mode selection, and PPO for hyperparameter control.
%
% 5. The contributions: A bulleted list of the paper's key contributions.
%    - Enhanced GVNS with 8 operators.
%    - RL-based adaptive operator selection (Q-learning vs DQN).
%    - ML-based FI/BI strategy selector (Random Forest).
%    - PPO-based hyperparameter controller.
%    - Extensive experiments on 50–400 customer instances.
%    - Open-source code and instances.
%
% 6. Paper structure: A brief roadmap of the sections.
%% ----------------------------------------------------------------------


% ----------------------------------------------------------------------
%% SECTION 2: PROBLEM DEFINITION AND BACKGROUND
% ----------------------------------------------------------------------
\section{Problem definition and Background}
\label{sec:background}

\subsection{VRPTW-S problem definition}
% Content to include (from report Chapter 1):
% - Formal definition of the Vehicle Routing Problem with Time Windows (VRPTW).
%    - Graph model G = (V, E), with vertex 0 as the depot.
%    - Definition of time windows [e_i, l_i], service times s_i, and travel costs c_ij.
% - Synchronization constraints: strict, overlap, precedence.
%    - Provide the Big-M formulation for a strict synchronization constraint.
%    - Briefly explain why this increases problem complexity.
% - The objective function: Minimizing total travel cost (or time).
% - Conclude by stating the NP-hard nature of the problem.

% Figures/tables to include:
% None (this is a conceptual sub-section).

\subsection{Related work}
% Content to include (from report Chapter 2, State of the art):
% - Exact methods (e.g., Bredström & Rönnqvist 2008) and their limitations.
% - Survey of metaheuristics: Simulated Annealing (Afifi et al. 2016), ALNS (Liu et al. 2019), VNS/GVNS (Frifita & Masmoudi 2018; Masmoudi et al. 2025).
% - Review of RL for routing:
%    - Early work on TSP/VRP (Bello et al. 2017; Kool et al. 2019).
%    - RL for operator selection (Xu et al. 2023).
% - Position your work: None have combined all three learning components for VRPTW-S.

\subsection{Reinforcement Learning preliminaries}
% Content to include (from report Chapter 3, Foundations of Machine Learning):
% - The Markov Decision Process (MDP) framework: formal definition of S, A, P, R, γ.
% - Value-based methods: Q-learning (tabular) and Deep Q-Network (DQN).
% - Policy-based methods: policy gradient and PPO.
% - The exploration vs exploitation dilemma (ε-greedy, UCB).

% Figures to include:
% - Figure 6 from your report (Agent-Environment interaction in RL).
%   Label as Fig. 1 in this paper.
\iffalse
\begin{figure}[htbp]
    \centering
    \includegraphics[width=0.7\linewidth]{path/to/agent-environment.pdf}
    \caption{The standard reinforcement learning interaction loop between an agent and its environment.}
    \label{fig:mdp_loop}
\end{figure}
\fi

\subsection{Machine Learning preliminaries}
% Content to include (from report Chapter 6, Learning the local search mode):
% - Brief explanation of bagging and feature randomness in Random Forests.
% - Distinction between a classifier (probability of improvement) and a regressor (expected gain).

% ----------------------------------------------------------------------
%% SECTION 3: BASELINE METAHEURISTIC – ENHANCED GVNS
% ----------------------------------------------------------------------
\section{Baseline metaheuristic: Enhanced GVNS}
\label{sec:gvns}

\subsection{GVNS principle}
% Content to include (from report Chapter 4, GVNS metaheuristic):
% - The core idea of VNS: a local optimum for one neighborhood is not necessarily a local optimum for another.
% - The GVNS architecture: Shaking phase → Variable Neighborhood Descent (VND) → Move acceptance.
% - The augmented evaluation function with dynamic penalties for time-window and horizon violations.

% Figures to include:
% - Conceptual diagram of the GVNS loop (combine Figs. 10-15 into a single flowchart).
%   Label as Fig. 2 in this paper.
\iffalse
\begin{figure}[htbp]
    \centering
    \includegraphics[width=0.8\linewidth]{path/to/gvns_flow.pdf}
    \caption{The principal architecture of the enhanced GVNS.}
    \label{fig:gvns}
\end{figure}
\fi

\subsection{Neighborhood operators}
% Content to include (from report Chapter 4, §4.3):
% - The original four operators: intra-route insertion, inter-route insertion, intra-route swap, inter-route swap.
% - The four newly added operators: cross exchange, k-relocate, 2-opt, 2-opt*.
% - For each, include a single sentence on its function (e.g., "The k-relocate operator moves a sequence of k consecutive visits...").
% - Why these additions were beneficial: they restore performance lost when re-implementing in Python.

% Table to include:
% - A new table summarizing the 8 operators, their type, and a brief description.
%   Label as Table 1.

\subsection{Variable Neighborhood Descent (VND)}
% Content to include (from report Chapter 4, §4.4):
% - Explanation of the fixed-order exploration of neighborhoods.
% - The distinction between First Improvement (FI) and Best Improvement (BI) strategies.
% - The pseudocode (as seen in Fig. 12 and Fig. 13 of your report).

% Figure to include:
% - Figure 12 (Pseudo-code of BVND) from your report, refined.
%   Label as Fig. 3.
% - Figure 13 (Pseudo-code of neighborhood exploration) from your report.
%   Label as Fig. 4.
\iffalse
\begin{algorithm}
\caption{Variable Neighborhood Descent (VND)}
\begin{algorithmic}
\STATE \textbf{Input:} Initial solution $S$, neighborhood set $\mathcal{N} = \{N_1, ..., N_{k_{max}}\}$
\STATE $k \leftarrow 1$
\WHILE{$k \le k_{max}$}
    \STATE $S' \leftarrow \text{Explore}(N_k, S)$
    \IF{$f(S') < f(S)$}
        \STATE $S \leftarrow S'$
        \STATE $k \leftarrow 1$ \COMMENT{Reset to first neighborhood}
    \ELSE
        \STATE $k \leftarrow k+1$ \COMMENT{Move to next neighborhood}
    \ENDIF
\ENDWHILE
\STATE \textbf{Output:} Improved solution $S$
\end{algorithmic}
\end{algorithm}
\fi

% ----------------------------------------------------------------------
%% SECTION 4: REINFORCEMENT LEARNING FOR OPERATOR SELECTION
% ----------------------------------------------------------------------
\section{RL for operator selection}
\label{sec:rl_operator}

\subsection{MDP formulation}
% Content to include (from report Chapter 5, Learning the operator order with RL):
% - How the operator selection problem is mapped to an MDP.
% - The State: The 24-dimensional feature vector (instance-level, solution quality, search progress, neighborhood success rates).
% - The Action Space: A = {N1,...,N8} representing the eight neighborhood operators.
% - The Reward Function: The normalized and shaped improvement.
% - The Transition Function: Deterministic, moving to the next state after the operator is applied.

% Figure to include:
% - Figure 18 (Global architecture of classical GVNS) and Figure 19 (RL-driven GVNS) from your report.
%   Combine or reference them to show the difference.
%   Label as Fig. 5.

\subsection{Tabular Q-Learning Agent}
% Content to include (from report Chapter 5, §5.3):
% - State discretisation into 144 abstract states (based on feasibility, stagnation, effectiveness, and instance size).
% - The Q-table structure and the Bellman update rule.
% - Action selection using the Upper Confidence Bound (UCB) algorithm for efficient exploration.
% - Key hyperparameters: learning rate α = 0.1, discount factor γ = 0.95, exploration constant c = 1.0.

% Table to include:
% - Table IX (Hyperparameters used for Q-learning agent) from your report.
%   Label as Table 2.

% Figure to include:
% - Figure 20 (Internal architecture of the Q-learning agent) from your report.
%   Label as Fig. 6.
\iffalse
\begin{figure}[htbp]
    \centering
    \includegraphics[width=0.8\linewidth]{path/to/ql_arch.pdf}
    \caption{The internal architecture of the tabular Q-learning agent.}
    \label{fig:ql}
\end{figure}
\fi

\subsection{Deep Q-Network (DQN) Agent}
% Content to include (from report Chapter 5, §5.4):
% - The need for function approximation due to the large state space.
% - The MLP architecture: Input (24 normalized features) → 2 x 128 hidden layers (ReLU) → Output (8 Q-values).
% - Stability mechanisms: Experience replay buffer (size 50,000), target network (updated every 100 steps), and reward clipping to [-1, 1].
% - Action selection using an epsilon-greedy policy (epsilon decays from 1.0 to 0.05 over 500 steps).
% - Key hyperparameters: learning rate 3x10^-4, batch size 32.

% Table to include:
% - Table X (Hyperparameters used for DQN agent) from your report.
%   Label as Table 3.

% Figure to include:
% - The MLP architecture figure from your report (labeled "Figure : Architecture of the neural network...").
%   Label as Fig. 7.

\subsection{Comparison of Q-Learning and DQN}
% Content to include (from report Chapter 5, synthesis):
% - A high-level conceptual comparison, positioning Q-learning as simple and effective for small state spaces, and DQN as robust and scalable for larger ones.
% - This section merely sets up the experimental comparison for the next part.

% ----------------------------------------------------------------------
%% SECTION 5: MACHINE LEARNING FOR LOCAL SEARCH MODE SELECTION
% ----------------------------------------------------------------------
\section{ML For local search mode selection}
\label{sec:ml_fibi}

\subsection{Problem statement: First vs. Best improvement}
% Content to include (from report Chapter 6, §6.1):
% - The trade-off: First Improvement (FI) is fast but may miss better moves.
% - Best Improvement (BI) is thorough but computationally expensive.
% - Neither strategy dominates; the best choice depends on the search context.

\subsection{Dataset collection and Feature engineering}
% Content to include (from report Chapter 6, §6.2):
% - How the dataset was collected: Running GVNS on training instances in both FI and BI modes.
% - Logging procedure: One sample per local search call, with features and outcomes.
% - The four categories of features (24 total): instance-level, dynamic search-level, solution-based structural, and neighborhood context.
% - The labeling and cleaning process: removing uninformative samples, creating improvement targets.

% Table to include:
% - Summary of the four feature categories (based on report's final retained features).
%   Label as Table 4.

\subsection{Two-Head Random Forest model}
% Content to include (from report Chapter 6, §6.3):
% - The Classifier Head: Trained to predict the probability of improvement, P(success).
% - The Regressor Head: Trained to predict the expected relative gain, E[gain].
% - The Combined Score: A decision is made by comparing Score_FI = P_FI × E[gain]_FI vs. Score_BI = P_BI × E[gain]_BI.
% - Model hyperparameters: number of trees = 400, max depth = 16, min samples per leaf = 30.

% Table to include:
% - Hyperparameters for the two-head RF model (from your report).
%   Label as Table 5.
% - Evaluation results for the classifier (ROC-AUC, Brier score) and regressor (R², MAE, RMSE).
%   Label as Table 6.

% Figure to include:
% - The internal layout figure of the RF model workflow (Figure --- from your report).
%   Label as Fig. 8.

\subsection{Integration into GVNS}
% Content to include (from report Chapter 6, §6.5):
% - The counterfactual evaluation process: The model predicts for both FI and BI simultaneously.
% - The fallback decision rule to avoid erratic behavior.
% - The pseudocode for the modified local search function.

% Figure to include:
% - The inference pipeline figure (Figure --- from your report).
%   Label as Fig. 9.

% Algorithm to include:
% - The pseudocode for the ML-guided BVND (Figure --- from your report).
%   Label as Algorithm 1.

% ----------------------------------------------------------------------
%% SECTION 6: PPO FOR HYPERPARAMETER CONTROL
% ----------------------------------------------------------------------
\section{PPO for hyperparameter control}
\label{sec:ppo}

\subsection{MDP formulation}
% Content to include (from report Chapter 7, The first sections):
% - Defining the problem as an MDP for controlling hyperparameters online.
% - The State: The current search context (from §6.2) plus the current values of the hyperparameters.
% - The Action Space: A set of 8 discrete actions, each triggering a specific behavioral change (e.g., Neutral, Diversify, Intensify, Reinforce Feasibility, Reset).
% - The Reward Function: The local improvement, normalized and compressed with a tanh function for stability.

% Table to include:
% - Table --- (Action space of the parameter controller) from your report.
%   Label as Table 7.

\subsection{PPO architecture and Offline training}
% Content to include (from report Chapter 7, §7.4 – §7.6):
% - The actor-critic architecture: networked with a shared trunk (two 128-neuron layers) for a compact design.
% - Explanation of the PPO clipped surrogate objective (ε=0.2) and value function loss.
% - The offline training procedure: dataset from GVNS runs, computing returns, advantages (GAE, λ=0.95), and batch updates.
% - The entropy bonus (coefficient 0.01) to encourage exploration.
% - Key hyperparameters (γ=0.99, learning rate 3x10^-4, batch size 256).

% Table to include:
% - Hyperparameters for PPO (from your report).
%   Label as Table 8.

% Figure to include:
% - The architecture of the PPO network (Figure --- from your report).
%   Label as Fig. 10.

\subsection{Integration into GVNS}
% Content to include (from report Chapter 7, §7.7):
% - How the controller is inserted into the main GVNS loop.
% - The concept of "fine granularity": The controller acts at every neighborhood evaluation.
% - The safety mechanisms: bounded and smoothed parameter updates.

% Figure to include:
% - The pseudo-code figure for PPO integration (Figure --- from your report).
%   Label as Fig. 11.

% ----------------------------------------------------------------------
%% SECTION 7: EXPERIMENTAL RESULTS
% ----------------------------------------------------------------------
\section{Experimental results}
\label{sec:results}

%% The following subsections will be filled from your Chaps. 5, 6, & 7.

\subsection{Experimental setup}
% Content to include (from report Chapter 4, §4.6):
% - The Solomon-based instances, sizes, and synchronization.
% - The baseline GVNS configuration.
% - The number of runs (5) and evaluation metrics (% gap, % improvement, etc.).

% \subsection{Instance generation for experimentation}
% Content to include (from report Chapter 4, §4.6):
% - Basis: Classical Solomon instances (C, R, RC distributions).
% - Instance sizes: 50, 100, 200, and 400 customers.
% - Procedure for adding synchronization constraints (5/10/20/40 synchronized pairs).
% - Number of vehicles: proportional to instance size (e.g., 25 vehicles for 50 customers).
% - Resulting dataset: 72 test instances.

\subsection{Results: Operator selection (Q-Learning vs. DQN)}
% Content to include (from report Chapter 5, §5.3.2 and §5.4.2):
% - A summary table comparing the performance of QL and DQN across all instance sizes.
%   This is a crucial combining of your two existing tables.
% - Analysis of QL's strengths (small instances) and failure (large instances).
% - Analysis of DQN's stability and scalability.

% Table to include:
% - A new master table comparing QL and DQN results.
%   Label as Table 9.

\subsection{Results: Local search mode selection (RF)}
% Content to include (from report Chapter 6, §6.7):
% - Performance of the RF model alone.
% - Performance of the hybrid models: QL+RF, DQN+RF.
% - Highlight the synergistic gains, especially for DQN+RF on large instances.

% Table to include:
% - A new table comparing the three FI/BI strategies.
%   Label as Table 10.

\subsection{Results: Hyperparameter control (PPO)}
% Content to include (from your report, to be completed):
% - Performance comparison of the static DoE configuration vs. the PPO-adaptive controller.
% - Report final cost gaps, runtime, and perhaps the parameter evolution over time.

% Table to include:
% - A new table comparing Static vs. PPO.
%   Label as Table 11.

% ----------------------------------------------------------------------
%% SECTION 8: OVERALL EVALUATION AND DISCUSSION
% ----------------------------------------------------------------------
\section{Overall evaluation and discussion}
\label{sec:evaluation}

\subsection{Performance summary}
% Content to include (synthesize from Chaps. 5, 6, & 7):
% - A holistic conclusion of the best-performing approach for different scenarios.
% - Clear recommendations: DQN-based operator selection is the foundation; for very large instances, add the RF model; the PPO controller is the next step for online adaptation.

% Table to include:
% - A master comparison table of ALL methods: GVNS, QL, DQN, RF, PPO, and hybrids.
%   Label as Table 12.

\subsection{Model behavior and Robustness analysis}
% Content to include:
% - A brief look at the decision behavior of the learned policies.
% - For example, mention how the policy's preference evolves (e.g., early preference for FI for speed, later preference for BI for exploitation).

\subsection{Limitations and Future perspectives}
% Content to include (from report Chapter 8):
% - Limitations: Hand-crafted state features, scalability beyond 400 customers, single-objective focus.
% - Future work: Graph Neural Networks for state representation, multi-objective RL, transfer learning.

% ----------------------------------------------------------------------
%% CONCLUSION
% ----------------------------------------------------------------------
\section{CONCLUSION}
\label{sec:conclusion}

% Content to include:
% - A powerful, concise restatement of the paper's core problem and its significance.
% - A summary of your proposed multi-part framework.
% - The key takeaway: Data-driven, adaptive search dramatically outperforms static optimization.
% - A final, forward-looking sentence on the impact of this research for operational research and logistics.


\section*{Acknowledgments}
% To be filled in later.

%% ----------------------------------------------------------------------
%% BIBLIOGRAPHY
%% ----------------------------------------------------------------------
%% Use a .bib file. This is a placeholder for the style.
\bibliographystyle{elsarticle-num-names}
\bibliography{your_bibliography_file}

\end{document}